\documentclass[10pt,a4paper]{article}
\usepackage[margin=1.6cm]{geometry}
\usepackage[T1]{fontenc}
\usepackage[utf8]{inputenc}
\usepackage{lmodern}
\usepackage{enumitem}
\usepackage{booktabs,tabularx,array}
\usepackage[hidelinks]{hyperref}
\usepackage[backend=biber,style=numeric,sorting=none,maxbibnames=3,doi=false,url=false,isbn=false,eprint=false,giveninits=true]{biblatex}
\addbibresource{refs.bib}
\setlist{nosep,leftmargin=1.2em,itemsep=1pt}
\setlength{\parindent}{0pt}
\setlength{\parskip}{3pt}
\usepackage[compact]{titlesec}
\titlespacing*{\section}{0pt}{8pt}{2pt}
\titlespacing*{\subsection}{0pt}{6pt}{1pt}
\usepackage{xcolor}
\renewcommand{\bibfont}{\scriptsize}

% Status of a question
\newcommand{\st}[2]{\colorbox{#1}{\footnotesize\strut #2}}
\newcommand{\ready}{\st{green!20}{ready to confirm}}
\newcommand{\leaning}{\st{yellow!35}{leaning}}
\newcommand{\open}{\st{red!15}{open}}
% Question block: ID(s), question, number of papers, status
\newcommand{\qblock}[1][]{\def\qkey{#1}\qblockx}
\newcommand{\qblockx}[4]{%
  \ifx\qkey\empty
  \subsection*{\underline{#1\quad #2}}%
  \else\hypertarget{\qkey}{}
  \subsection*{#1\quad #2}\fi%
{\footnotesize #3 papers\quad}#4\par}
\newcommand{\more}[2]{\hyperlink{#1}{\colorbox{blue!10}{\footnotesize\strut\textcolor{blue!60!black}{\underline{#2}\,$\rightarrow$}}}}
\newcommand{\sitem}[5]{
\item \textbf{#1}\par\vspace{1pt}
  \begin{itemize}[leftmargin=1.1em,itemsep=1pt,label={\textbullet}]
    \item #2
    \item \textit{Lean:} #3
    \item \textit{To discuss:} #4
    \item \textit{Evidence:} #5
\end{itemize}}
\newcommand{\ev}{\textbf{Evidence.}}
\newcommand{\lean}{\textbf{Current lean.}}
\newcommand{\ask}{\textbf{To discuss.}}

\begin{document}

\begin{center}
  {\large\bfseries Open questions and the evidence so far: supervisor meeting
  2026-11-24}\\[2pt]
  % {\small Companion to the literature summary. Each block
  %   combines what the papers say about one thesis question (or a group of
  %   closely related ones), then gives the current lean and what to discuss.
  %   Dataset questions come first (deadline: dataset frozen 2027-01-22), then
  %   agent, evaluation and framing. Within each part, questions with the most
  %   evidence come first. ``Papers'' counts the reviewed papers cited in the
  %   block. Numbers are from the papers (section/table in brackets).
  %   \ready\ enough evidence to decide now,
  %   \leaning\ a preferred option exists but is not decided,
  % \open\ no preferred option yet.}
\end{center}

\section*{Summary for the meeting}
{\small
  \begin{enumerate}[leftmargin=1.5em,itemsep=7pt]
      \sitem{Is there a gap for a new dataset?}
      {Wikidata is widely used to \emph{build} visual QA data, but no
        benchmark releases image + question + executable SPARQL, and
      none evaluates an agent that queries the graph.}
      {The gap holds, so the dataset can be presented as a contribution.}
      {Is that enough, given that the research questions are only
      about the agent?}
      {\more{lit}{Existing benchmarks}}
      \sitem{How is difficulty defined?}
      {Hops, query complexity, image role (subject only, a related
        entity, or a busy scene, which lowers retrieval a lot) and
      entity popularity all affect accuracy.}
      {Levels set a priori from hops, query complexity and image
        role; popularity and question type logged as non-exclusive
        labels; hop cap of 3 (4 only for a few hard items); checked on
      a pilot run before the freeze.}
      {A priori levels, from pilot results, or both?}
      {\more{diff}{Difficulty} \more{types}{Types} \more{image}{Image
      role} \more{hops}{Hops}}
      \sitem{How are items curated and accepted?}
      {Questions are suggested by an LLM and verified by hand
      (precedents exist). Unreviewed data is often ambiguous or wrong.}
      {Accept an item only if the answer is unique and supported, the
        entity is not named, the image is needed and the gold query
        runs. The gold query returns only the answer column as
      identifiers; its result is frozen with a date.}
      {Keep yes/no questions or rephrase them as SELECT? One or
      several gold queries per item?}
      {\more{cur}{Curation} \more{gold}{Gold queries}}
      \sitem{Which distractors go into the shared image index, and how many?}
      {Most benchmarks use all other entities; similar distractors
      and a larger index clearly hurt retrieval.}
      {Same-class distractors (same \texttt{instance of}), optionally
        plus text-similar neighbours. No preference on the count yet;
      it could be an experiment axis.}
      {Same class only, or mixed?}
      {\more{dist}{Distractors}}
      \sitem{How large should the dataset be?}
      {Curated SPARQL benchmarks have roughly 100 to 500 questions.}
      {Target about 300, curated manually before the freeze on
        2027-01-22. Also depends on how many examples could we adapt
      and reuse from the VQA datasets.}
      {Is that realistic? Is reporting diversity worthwhile?}
      {\more{size}{Dataset size}}
      \sitem{Which agent and tools?}
      {Search tools matter most. GRASP is open and strong, so it is a
      possible external baseline once given an image tool.}
      {Own agent with the Wikidata search and query tools of SPINACH,
        the constrained searches of GRASP, and the image-index search
        as one more tool. A crop tool is liked (improves grounding,
      possible ablation) but not decided.}
      {Is a crop tool in scope? Is GRASP with an image tool needed,
        or are the internal configurations plus a direct image +
      question VLM answer enough?}
      {\more{tools}{Tools} \more{base}{Baselines}}
      \sitem{Configurations, metrics and ablations}
      {Telling the model the entity helps in every paper read.
      Row-wise F1 / exact match is the de-facto metric for SPARQL results.}
      {Five configurations: oracle with the entity name, oracle with
        its QID, vision only, index only, both. Primary metric row-wise
        F1 and exact match; secondary grounding accuracy, cost,
        iterations. Ablations (distractor count, tool removal, crop
      variants, error analysis) limited by model capacity.}
      {Add a question-only control? Is a label-based metric needed to
      compare with the visual papers?}
      {\more{cfg}{Configurations} \more{metrics}{Metrics}
      \more{ablat}{Ablations}}
      \sitem{Research questions and hypotheses}
      {Oracle gaps support the hypothesis that most of the gap is
        grounding. That the index beats vision for less known entities
      matches the recognition literature, but not the answer-accuracy results.}
      {State the hypothesis on grounding accuracy.}
      {Is that the right measure?}
      {\more{rq}{Research questions}}
  \end{enumerate}
}
\newpage

\section{Dataset and image index}

\qblock[lit]{LIT-01}{Does an image + question + gold-SPARQL benchmark
already exist?}{11}{\ready}
\ev
\begin{itemize}
  \item Wikidata is often used to \emph{build} visual QA data, but the SPARQL
    is never released. KVQA withholds its queries on purpose
    \cite[S3]{shah2019kvqa}; M$^3$-VQA uses SPARQL only while building
    questions and releases neither SPARQL nor QIDs \cite[\S\,4.1]{ma2026m3};
    ReasonVQA releases a QID route and the last property, no query
    \cite[App.~1.2]{tran2025reasonvqa}; InfoSeek turns single Wikidata triples
    into templates \cite[\S\,3.2]{chen2023can}.
  \item The other visual benchmarks use Wikipedia text as the knowledge base
    (Encyclopedic-VQA, RETINA, the KB-VQA audit)
    \cite{mensink2023encyclopedic,lee2025breaking,ma2026identifying}, or stop
    at entity recognition (OVEN, KnowCoL) \cite{hu2023open,zhou2025seeing}.
  \item None evaluates an agent that queries a KG; the only agent
    (M$^3$-VQA) retrieves Wikipedia text and images \cite[\S\,5.4]{ma2026m3}.
    Text-only Wikidata benchmarks have gold SPARQL but no images
    \cite{loustaunau2025qawiki,walter2026grasp}.
\end{itemize}
\lean\ The gap holds: (image, entity-free question, executable gold SPARQL)
+ a shared image index with distractors + an agentic text-to-SPARQL
evaluation. Still unchecked: WikiVQABench (2026, multiple choice grounded in
Wikidata triples) and ViQuAE.\\
\ask\ Is this gap enough to present the dataset as a contribution
(D-017), given that the RQs are only about the agent?

\qblock[diff]{DS-13}{How is difficulty defined and measured?}{9}{\open}
\ev\ Every axis below lowers accuracy in at least one paper:
\begin{itemize}
  \item \emph{Hops.} M$^3$-VQA accuracy falls with hops and entities
    (\S\,5.5) \cite{ma2026m3}; ReasonVQA: 3-hop clearly harder, 1-hop barely
    easier than 2-hop (\S\,4.4) \cite{tran2025reasonvqa}; E-VQA two-hop 22.8
    vs.\ single-hop 48.8 (\S\,5.6) \cite{mensink2023encyclopedic}.
  \item \emph{Query complexity.} KGVoyager counts triple patterns and
    special constructs; its lead over GRASP shrinks at 7+ triple patterns
    (App.~C) \cite{wisam2026kgvoyager}.
  \item \emph{Grounding.} Image shows a related entity: Recall@5 51.7 / 64.9
    $\to$ 20.8 (RETINA, Table 1) \cite{lee2025breaking}; one extra entity in
    the image: R@1 43.5 $\to$ 14.7 / 20.4 (Table 2) \cite{ma2026identifying};
    more objects in the image: 41.6 $\to$ 37.1 (Table 3)
    \cite{jian2024large}; distractors, see IX-01.
  \item \emph{Popularity} changes which approach wins. For recognition, OVEN:
    PaLI wins on head entities ($>$5K monthly views), CLIP2CLIP is on par or
    better on tail ones ($<$2.5K) (Fig.~6) \cite{hu2023open}. For answer
    accuracy, InfoSeek points the other way, but its pipeline also needs a
    text QA model to know facts about the entity (Fig.~7) \cite{chen2023can}.
  \item \emph{Answer type.} Time and numerical answers are hardest in InfoSeek
    (Table 6) \cite{chen2023can}, but it's based on text passages
    rather than properties in Wikidata.
\end{itemize}
\lean\ Define difficulty a priori from properties of the question
(hops, query complexity, image role), log popularity as a label, and use
the pilot (2026-12-07 to 12-18), including an image-only run without the
question text, to check the spread before the freeze.\\
\ask\ A priori levels, a posteriori from pilot results, or both? Which
two or three axes should define easy / medium / hard?

\qblock[cur]{DS-09, DS-12}{How are questions curated, and what are the
acceptance criteria?}{8}{\ready}
\ev
\begin{itemize}
  \item LLM-suggested, human-verified construction has precedents: SPINACH
    (GPT-4o proposes, experts verify by running intermediate queries,
    App.~C.3) \cite{liu2024spinach} and QAWiki (LLM phrasing suggestions
    reviewed by hand, plus 14 quality-control queries, \S\,3.7)
    \cite{loustaunau2025qawiki}. This supports D-016.
  \item Unreviewed data is often wrong or ambiguous: 22\,\% of InfoSeek
    answers are not supported by the evidence, 47\,\% / 59\,\% of InfoSeek /
    E-VQA questions are ambiguous \cite[\S\,3.1--3.2]{ma2026identifying};
    OVEN finds 8.4\,\% wrong entity links \cite[App.~10.2]{hu2023open}; RETINA
    keeps only 29.2\,\% of 6K generated questions after curation, mostly
    because several answers were valid \cite[App.~A.3]{lee2025breaking}.
  \item The review checks in these papers reduce to the same few:
    \emph{answer supported} (InfoSeek consistency filter, the audit),
    \emph{unique answer} (the audit, RETINA's qualifying entity, OVEN's
    rewriting), \emph{no leakage} (demonstrative instead of the name,
    BM25 check, removal of questions solvable without the image; RETINA),
    and \emph{no quality issue} (ambiguous, subjective, incomplete, ties;
    QAWiki \S\,3.5)
    \cite{chen2023can,ma2026identifying,lee2025breaking,hu2023open,loustaunau2025qawiki}.
  \item Query images: RETINA keeps query and document images in separate
    pools with a perceptual-hash check (hash size 8, distance $\le 10$,
    App.~A) \cite{lee2025breaking}; D-005 already requires a different
    Commons photo than the index image.
  \item Template pipelines (InfoSeek, ReasonVQA, M$^3$-VQA) give scale but
    no released gold query. ReasonVQA items could seed ideas (QID routes,
    landmarks only, AGPL) \cite{tran2025reasonvqa}; paraphrasing as in KVQA
    costs 17 points (49.5 $\to$ 32.0, Table 4) \cite{shah2019kvqa}.
\end{itemize}
\lean\ Accept an item only if: one correct, unambiguous answer; no
quality-tag issue (exclude instead of tagging, unlike QAWiki); answer stable
in the snapshot; entity not named and no other leakage; the image is needed;
the gold SPARQL runs and its result supports the answer; natural phrasing;
image licence allows reuse. BM25 and perceptual-hash checks automated, the
rest manual.\\
\ask\ Confirm the list. Can ReasonVQA / RETINA test items be used as
seeds for ideas (each still needs our own gold query)?

\qblock[dist]{IX-01, IX-02, LIT-04}{Which distractors go into the index,
and how many?}{8}{\leaning}
\ev
\begin{itemize}
  \item Almost nobody selects distractors: KVQA, OVEN and InfoSeek use
    every other entity of the KB (InfoSeek adds frequent Wikipedia entities)
    \cite{shah2019kvqa,hu2023open,chen2023can}; M$^3$-VQA samples 5 of the
    20 nearest neighbours by \emph{text} summary (App.~D.3) \cite{ma2026m3};
    ReasonVQA only has random multiple-choice answers \cite{tran2025reasonvqa}.
  \item Visually similar distractors hurt: in RETINA, documents whose
    main entity matches
    the image are retrieved more often than the gold one (Fig.~5)
    \cite{lee2025breaking}; in PMSR, even with the oracle text, 10 retrieved
    image-text pairs lower accuracy 86.6 $\to$ 78.7\,\%, 20 pairs to
    64.3\,\% (Table A8) \cite{choi2025progressive}; in the POC the egg chair
    was not in the top 5 among 25 same-class entities.
  \item Size matters: OVEN CLIP2CLIP falls roughly linearly in log(candidates)
    from 20K to 6M (Fig.~5) \cite{hu2023open}; KVQA loses over 10 points from
    18K to 69K identities \cite{shah2019kvqa}.
  \item The KB-VQA audit puts the distractor into the query image, so it is no
    evidence for the index \cite{ma2026identifying}.
\end{itemize}
\lean\ Same-class distractors (\texttt{instance of}) in the one shared
index (D-006); text-similar neighbours or a frequent-entity pool as
additions. No preference on the count yet.\\
\ask\ Same class only, or mixed? Should the count be varied as an
experiment axis (RQ3)?

\qblock[types]{DS-01, DS-14, LIT-02}{Which question types, and which labels
per item?}{7}{\leaning}
\ev
\begin{itemize}
  \item The papers classify along different axes: answer type (InfoSeek
      time / numerical / string; QAWiki entity, entity set, numeric, temporal,
    boolean, other) \cite{chen2023can,loustaunau2025qawiki}; structure (KVQA's
      ten non-exclusive categories such as Boolean, comparison, counting,
    multi-entity; M$^3$-VQA parallel vs.\ sequential)
    \cite{shah2019kvqa,ma2026m3}; hops and two-hop chaining through a bridge
    entity \cite{tran2025reasonvqa,mensink2023encyclopedic}; knowledge domain
    by property (ReasonVQA, 20 domains).
  \item Non-exclusive tags allow results per category: KVQA oracle accuracy
    ranges from Boolean 75.1 to subtraction 40.5 (Table 5) \cite{shah2019kvqa}.
    Template datasets (ReasonVQA) have no comparison, count or Boolean
    questions.
  \item Popularity (page views) is split on in InfoSeek and OVEN
    \cite{chen2023can,hu2023open}.
\end{itemize}
\lean\ Several independent label axes per item: hop count, structural
tags (starting from KVQA's categories), answer type (QAWiki), and as
candidates popularity, image role and class of the depicted entity.
QAWiki's property-mention types: too costly. KGQA taxonomies (Mintaka,
LC-QuAD 2.0, ComplexWebQuestions) not read yet.\\
\ask\ Non-exclusive tags or one type per question? Quotas per type? Comparison
questions need a second entity named in the text (multi-image is a stretch
goal); acceptable?

\qblock{IX-03, IX-04}{Which embedding model, and what does the index
store?}{6}{\leaning}
\ev
\begin{itemize}
  \item Every paper represents an entity with image \emph{and} text:
    M$^3$-VQA sums CLIP embeddings of summary and image, better than either
    alone (App.~F.4) \cite{ma2026m3}; OVEN name + first infobox image
    \cite{hu2023open}; KnowCoL description + mean of lead images
    \cite{zhou2025seeing}; PMSR image + text beats image only, R@5 43.1 vs.\
    36.2 (Table A1) \cite{choi2025progressive}.
  \item A frozen CLIP index without training is used by SAGE-VQA
    \cite[App.~A.7]{shen2026sage}; LLM-RA trains its retriever
    \cite{jian2024large}. Larger embedding models help a little (PMSR App.~J).
\end{itemize}
\lean\ CLIP ViT-B/32 as in the POC; payload QID only (the agent looks up
the rest with SPARQL); other models only if time allows.\\
\ask\ The literature mixes in text; keep image-only retrieval (cleaner
grounding test) or add a text variant as an ablation?

\qblock[image]{DS-05, DS-07}{What role does the image play?}{5}{\leaning}
\ev
\begin{itemize}
  \item Related-entity images (RETINA) make retrieval much harder (Table 1)
    \cite{lee2025breaking}; for Wikidata the related entity would be a graph
    neighbour of the subject, and the gold query follows from the graph.
  \item Busier images are harder: M$^3$-VQA has several entities per image
    \cite{ma2026m3}; ReasonVQA refers to one object among others via
    scene-graph context and these questions score lower (\S\,4.4)
    \cite{tran2025reasonvqa}; LLM-RA recall falls with object count
    \cite{jian2024large}; the audit composes a second entity into the image
    \cite{ma2026identifying}.
\end{itemize}
\lean\ Default: the image shows the subject. RETINA-style related-entity
images as an option for harder items (increasingly likely); prefer real
multi-entity images over automatic composites.\\
\ask\ Related-entity images in the first version or as an extension? Is
the ``object in the person's right hand'' idea still wanted?

\qblock{DS-04}{Which endpoint and snapshot?}{5}{\open}
\ev
\begin{itemize}
  \item Live Query Service: SPINACH (60\,s limit) and QAWiki, which lists
    changing results as a limitation (\S\,4.4)
    \cite{liu2024spinach,loustaunau2025qawiki}.
  \item Fixed version: KVQA uses a 2018 dump \cite{shah2019kvqa}; WikiKGQA
    2026 fixes one Wikidata version, and the GRASP authors load it into
    their own QLever for control over load and limits (\S\,2--3)
    \cite{walter2026grasp}. GRASP uses the public QLever endpoints and builds
    its own search index (6\,GB) \cite{walter2025grasp}.
\end{itemize}
\lean\ No lean on the endpoint; gold results are stored with a date in any
case. For entity search, the Wikidata search API (as SPINACH) is probably
enough.\\
\ask\ Is a dated gold result enough, or should agent runs also hit a fixed
snapshot (own QLever on a dump)?

\qblock[gold]{DS-08}{How are gold queries and gold results defined and
verified?}{4}{\leaning}
\ev
\begin{itemize}
  \item Answer form is a big error source: 18\,\% of KGVoyager's complete
    failures are representation errors (extra label columns, SELECT instead
    of ASK, label instead of IRI; App.~F) \cite{wisam2026kgvoyager}.
  \item The standard F1 ignores extra predicted columns
    \cite[\S\,5.1]{liu2024spinach}, so the gold query should project only the
    answer column(s), as identifiers (QAWiki, \S\,4.1)
    \cite{loustaunau2025qawiki}. Scoring only results can credit a wrong ASK
    query (QAWiki \S\,4.4).
  \item Yes/no questions need an ASK query, which returns a single true or
    false. A wrong ASK query that happens to return the right value gets full
    credit, and with a boolean answer a guess matches 50\,\% of the time, so
    the score says little about whether the image was grounded correctly.
    A SELECT query is much harder to get right by accident.
  \item Verification: SPINACH keeps gold queries that are non-empty and under
    10\,MB \cite{liu2024spinach}; the audit's check that the answer follows
    from the evidence means running the gold query on the snapshot for us
    \cite{ma2026identifying}.
  \item QAWiki writes queries as generally as possible and allows several
    gold queries per question (\S\,3.4).
\end{itemize}
\lean\ Gold query hard-codes the depicted entity (\texttt{BIND(wd:Q\ldots)});
answer columns only, as IRIs; result frozen with a date; items with ties or
incomplete data excluded. Unsure about several gold queries per item.\\
\ask\ Keep yes/no (ASK) questions or rephrase them as SELECT? One or
several gold queries?

\qblock[size]{DS-10}{How large should the dataset be?}{4}{\open}
\ev\ Curated KGQA sets are small: SPINACH 155 dev + 165 test, where the dev
set separates systems differing by at least 9.0\,\% EM (App.~C.2)
\cite{liu2024spinach}; KGVoyager's benchmarks 88--108 questions
\cite{wisam2026kgvoyager}; QAWiki 526 \cite{loustaunau2025qawiki}.
Diversity can be reported for a small set as mean pairwise embedding distance
(E-VQA, Table 3) \cite{mensink2023encyclopedic}.\\
\ask\ Is SPINACH scale (about 300) a sensible target with manual curation
by 2027-01-22? Report diversity?

\qblock[hops]{DS-03}{What is the maximum hop depth?}{3}{\open}
\ev\ ReasonVQA caps at 3 hops because longer generated questions became
ungrammatical, and only 5.4K of 4.2M are 3-hop \cite{tran2025reasonvqa};
M$^3$-VQA goes to 4+ (32.8\,\% of items) \cite{ma2026m3}; E-VQA stops at 2
\cite{mensink2023encyclopedic}. The POC has 2--3 hops, and recognising the
entity is an extra step before the first hop.\\
\ask\ Cap at 3, or allow a few 4-hop items for the hard end?

\section{Agent, evaluation and framing}

\qblock[tools]{AG-02}{Which tools does the agent get?}{9}{\leaning}
\ev
\begin{itemize}
  \item Search tools matter most: SPINACH dev F1 46.4 falls to 25.3 without
    search, 29.4 without property examples, 36.4 without the entity view
    (Table 4) \cite{liu2024spinach}; in GRASP execute-only is worst and
    constrained searches (properties of an entity, objects of a property)
    beat more flexible ones (Table 6) \cite{walter2025grasp}; GRISP loses 3.7
    F1 without this kind of KG guidance (Table 5) \cite{walter2026grisp};
    KGVoyager's hybrid search beats keyword-only by 1.1--26.5 F1 and drops
    low-scoring results to limit context length (Table 4)
    \cite{wisam2026kgvoyager}.
  \item Tool output is capped: first and last 5 rows, top 10 search results
    \cite{liu2024spinach,walter2025grasp}.
  \item Image side: cropping the key entity before retrieval raises InfoSeek
    Recall@5 39.8 $\to$ 47.3 (Table 4) \cite{jian2024large}; M$^3$-VQA's
    agent detects and crops \cite{ma2026m3}; SAGE-VQA checks each crop with
    the VLM first, but never measures the check alone \cite{shen2026sage}.
    PMSR's retrieve--condense--requery loop is an optional grounding mode
    \cite{choi2025progressive}.
  \item Exploration notes and a class shape index help GRASP unevenly
    ($-$2.7 to +5 F1) \cite{walter2026grasp}.
\end{itemize}
\lean\ SPINACH's Wikidata tools plus GRASP's constrained searches, and the
image-index search as one more tool in the same loop; high-level zero-shot
prompt. Crop tool (one merged check of crop and whole image) liked, not
decided; it needs image input (only Gemma 4 E4B on CampusAI).\\
\ask\ Is a crop tool in scope? Does the image tool take a handle or pixels?

\qblock[metrics]{LIT-03, EV-01}{Which metrics?}{9}{\ready}
\ev
\begin{itemize}
  \item One de-facto standard: SPINACH's row-wise F1 (optimal row matching,
    extra predicted columns not penalised, EM iff F1 = 1)
    \cite[\S\,5.1]{liu2024spinach}, reused by GRASP, KGVoyager, QAWiki and
    GRISP
    \cite{walter2025grasp,wisam2026kgvoyager,loustaunau2025qawiki,walter2026grisp}.
  \item Known weaknesses: values in the wrong column still count, wrong ASK
    queries can be credited \cite{loustaunau2025qawiki}; ambiguous questions
    and different LIMITs score unfairly low \cite{walter2025grasp}.
  \item Alternatives: QALD-F1 (WikiKGQA) gives abstaining P = 1 and scores
    higher \cite{walter2026grasp}; the visual papers score answer text (set
    IoU with aliases, exact / substring / semantic match, BEM)
    \cite{ma2026m3,tran2025reasonvqa,mensink2023encyclopedic}.
  \item Secondary measures: cost, iterations, time, tool calls
    \cite{wisam2026kgvoyager}; runtime \cite{walter2026grisp}.
\end{itemize}
\lean\ Primary: row-wise F1 and EM. Secondary: grounding accuracy (gold
QID in the index top-$k$; gold QID in the final query), cost, iterations,
tool use. QALD-F1 not planned.\\
\ask\ Is a label-based answer metric needed for comparison with the
visual papers?

\qblock[cfg]{AG-05}{Which agent configurations?}{8}{\ready}
\ev\ Telling the model the entity helps in every paper: InfoSeek
$\sim$20 $\to$ $\sim$50 \cite{chen2023can}, M$^3$-VQA 24.5 $\to$ 34.7
\cite{ma2026m3}, KVQA 49.5 $\to$ 59.2 \cite{shah2019kvqa}, E-VQA 19.7 $\to$
31.0 \cite{mensink2023encyclopedic}, and on text-only Wikidata QA, giving
the linked QIDs lifts F1 0.164 $\to$ 0.406 \cite{loustaunau2025qawiki}.
Image + question without tools is ReasonVQA's whole setting (GPT-4o 62.8)
\cite{tran2025reasonvqa}. Crop variants fit better as ablations
\cite{jian2024large,shen2026sage}.\\
\lean\ Five configurations: oracle with the name, oracle with the QID
(decided after QAWiki), vision only, index only, both. Image-only check
(no question text) in the pilot. An oracle that hands over triples is
unlikely.\\
\ask\ Also a question-only control (no image), as in M$^3$-VQA (15.6 vs.\
24.5) and InfoSeek?

\qblock[ablat]{EV-03}{Which ablations?}{7}{\open}
\ev\ Templates: remove one tool at a time (SPINACH Table 4, KGVoyager
Table 4) \cite{liu2024spinach,wisam2026kgvoyager}; function sets, model
size, a final feedback step (GRASP \S\,4.3) \cite{walter2025grasp}; switch
components off (PMSR Table 6) \cite{choi2025progressive}; whole image vs.\
crops vs.\ checked crops, only if a crop tool exists
\cite{jian2024large,shen2026sage}; hand-labelled error analysis of 20
(SPINACH) or 100 (OVEN) cases \cite{hu2023open}. KGVoyager runs each setting
3 times and some effects are within run-to-run noise.\\
\ask\ With limited CampusAI capacity, which two or three: distractor
count, tool removal, crop variants, error analysis? Repeated runs needed?

\qblock[base]{AG-01, EV-02}{Own agent or GRASP; which baselines?}{6}{\leaning}
\ev\ GRASP is open (Apache-2.0), wins WikiKGQA 2026 (85.0) and beats the
fine-tuned GRISP on QALD-7, QALD-10 and SPINACH
\cite{walter2025grasp,walter2026grasp,walter2026grisp}. GRISP needs
training pairs, so it is no baseline here (no train split, D-008).
KGVoyager does not port to Wikidata (no \texttt{rdf:type})
\cite{wisam2026kgvoyager}. A Claude Code agent matches GRASP, but paid APIs
are out of scope (D-010). PMSR is a fixed workflow, not ReAct
\cite{choi2025progressive}. A direct image + question VLM answer is a
ready-made baseline (ReasonVQA) \cite{tran2025reasonvqa}.\\
\lean\ Own agent, with GRASP and KGVoyager tool specifications as
inspiration; GRASP (plus an image tool) as a possible external baseline.\\
\ask\ Is GRASP with an image tool needed, or are the internal
configurations plus a VLM direct answer enough?

\qblock[rq]{FR-01}{Research questions and hypotheses}{5}{\leaning}
\ev\ The oracle gaps above support the first hypothesis (the agent
  falls clearly below the oracle,
  and most of the gap is failing to identify the entity rather than writing
wrong SPARQL) \cite{loustaunau2025qawiki,ma2026m3}. The second hypothesis
(the image index beats the model's own vision for less known entities, vision
beats it for well-known ones, the combination is best) matches OVEN's
recognition result
\cite{hu2023open}; InfoSeek's opposite result is for answer accuracy with a
Wikipedia-text pipeline \cite{chen2023can}. The GRASP authors name Wikidata
in LLM training data as a threat to validity (\S\,5) \cite{walter2026grasp}.\\
\lean\ State the index-versus-vision hypothesis on grounding accuracy
(OVEN), and mention InfoSeek as a
caveat for final answers.\\
\ask\ Is grounding accuracy the right measure for the
index-versus-vision hypothesis?

\qblock{AG-04}{Which CampusAI model?}{4}{\leaning}
\ev\ GRASP with open models: Qwen2.5-72B 68.7 vs.\ GPT-4.1 75.3; Qwen3-32B
61.3; a sharp drop below about 14B / 8B (Tables 5, 7) \cite{walter2025grasp}.
The backbone moves GRASP by about 9 F1 points \cite{walter2026grasp}. Smaller
models tend to give up in long contexts (50 of 56 failures, App.~F)
\cite{wisam2026kgvoyager}. A crop check needs an extra image call per crop
\cite{shen2026sage}.\\
\lean\ Qwen3.8-27B (tools, 224K context) is the obvious agent candidate;
Gemma 4 E4B is the only model with image input; cap tool output.
CampusAI contact not asked yet (rate limits,
pinned weights).

\section*{Not covered in the summary}
Underlined titles above (embedding model and index content, endpoint and
snapshot, CampusAI model) have evidence but are left out of the summary. The
following have little or no literature evidence yet:
\begin{itemize}
  \item DS-02: which part of Wikidata
  \item DS-06: multi-image questions (stretch)
  \item DS-11: dev/test split
  \item LIT-05: visual knowledge graphs (only ReasonVQA's scene graphs so far)
  \item AG-03: agent framework
  \item RP-01 to RP-06: repository
\end{itemize}

\printbibliography
\end{document}
